\documentclass[12pt,a4paper]{article}

\usepackage[utf8]{inputenc}
\usepackage[T1]{fontenc}
\usepackage[brazilian]{babel}
\usepackage{lmodern}
\usepackage[a4paper,top=3cm,bottom=2cm,left=3cm,right=2cm]{geometry}
\usepackage{setspace}
\usepackage{booktabs}
\usepackage{float}
\usepackage{tabularx}
\usepackage{enumitem}
\usepackage{url}
\usepackage[hidelinks]{hyperref}
\def\UrlBreaks{\do\/\do-\do_\do.\do a\do e\do i\do o\do u}

\onehalfspacing
\setlist{itemsep=2pt,topsep=4pt}

\title{Guia de Requerimentos UCSal\\[4pt]
\large Proposta de solução --- Programação Orientada a Objetos Avançada}
\author{João Vitor Bomfim de Souza\\
\small Engenharia de Software --- Universidade Católica do Salvador}
\date{Outubro de 2026}

\begin{document}

\maketitle

\section{Problema}

Quem já precisou abrir um requerimento no Portal do Aluno da UCSal sabe que a parte difícil não é
acompanhar o pedido depois de aberto. O portal (TOTVS Educacional) mostra número de protocolo, status
e o histórico de etapas por setor. O que falta é informação \emph{antes} de clicar em ``Solicitar''.

Fomos ao portal com uma conta de aluno ativa, em outubro de 2026, e registramos o que aparece em
\emph{Solicitações $\rightarrow$ Requerimentos}:

\begin{itemize}
  \item a aba ``Disponíveis'' lista oito tipos de requerimento. Todo o resto (declarações, histórico,
        trancamento, transferência externa, FIES, colação de grau) fica escondido dentro de um tipo
        genérico chamado ``Abertura de Requerimento'', em uma lista suspensa com 37 categorias
        identificadas apenas pelo nome;
  \item ao abrir os detalhes desse tipo genérico, os campos \emph{Descrição} e \emph{Procedimentos}
        estão vazios;
  \item no formulário, a coluna \emph{Prazo} da tabela de entrega também está vazia. Não há lista de
        documentos exigidos nem indicação de quais setores vão analisar o pedido. O aluno escreve o
        que acha que deve escrever em um campo de texto livre.
\end{itemize}

Três requerimentos reais de um integrante da equipe mostram onde isso pesa:

\begin{table}[H]
\centering
\small
\begin{tabularx}{\textwidth}{@{}lccX@{}}
\toprule
Requerimento & Aberto em & Dias até fechar & O que aconteceu \\
\midrule
Dispensa de estágio & 03/08/2025 & 32 &
  Indeferido por ``campo equivocado''. A categoria escolhida servia só para o curso de Direito,
  mas não existe outra para encerrar um estágio. A resposta não indicou o caminho certo. \\
\addlinespace
Trancamento de disciplina & 28/08/2026 & 31 &
  O Registro Acadêmico resolveu em 3 dias. Depois o pedido ficou 24 dias no Financeiro para
  recalcular parcelas, etapa que o aluno não sabia que existia. \\
\addlinespace
Transferência interna & 27/06/2025 & 10 &
  Deferido, mas terminou com ``ajuste as disciplinas se necessário'', devolvendo trabalho ao aluno. \\
\bottomrule
\end{tabularx}
\caption{Protocolos de um integrante da equipe no Portal do Aluno (dados consultados em 21/10/2026).}
\end{table}

O primeiro caso é o mais claro. O aluno tinha urgência, porque um novo estágio dependia da baixa do
anterior. Escolheu a categoria que parecia mais próxima, esperou um mês e recebeu uma negativa por ter
usado o campo errado. O erro não era de má vontade de ninguém: o próprio sistema não diz para que serve
cada categoria.

Isso não é um caso isolado. No Reclame Aqui, a UCSal tem 69 reclamações registradas e nota média 6,16
no período de março a agosto de 2026. Os assuntos que mais se repetem são acesso ao portal e ao AVA,
requerimentos sem retorno (trancamento, aproveitamento de disciplinas, reembolso) e atendimento por
WhatsApp que para de responder. Em uma delas, o aluno diz que não conseguia abrir nenhum requerimento
para pedir dispensa de disciplinas; em outra, que o sistema novo tirou a ferramenta de acompanhamento
que existia antes.

Resumindo: o aluno abre o pedido no escuro. Quando erra, só descobre semanas depois, e a secretaria
gasta tempo analisando pedidos que nunca poderiam ser deferidos naquele campo.

\section{Público-alvo}

\paragraph{Usuário principal.} Alunos de graduação da UCSal, presencial e EAD, que precisam abrir algum
requerimento. Alguns grupos sentem mais o problema:
\begin{itemize}
  \item calouros e alunos transferidos, que ainda não sabem como a universidade funciona por dentro;
  \item bolsistas ProUni e alunos com FIES, porque quase todo pedido deles passa também pela análise
        da bolsa ou do Financeiro;
  \item alunos que trabalham, que dependem de declarações e de trancamentos com prazo curto.
\end{itemize}

\paragraph{Beneficiados indiretos.} O Registro Acadêmico e a secretaria, que passariam a receber menos
pedidos na categoria errada, e o DCE ou centros acadêmicos, que hoje respondem essas dúvidas um a um em
grupos de WhatsApp.

\section{Solução}

\subsection{Ideia central}

A primeira tentação seria construir um sistema de chamados paralelo ao portal, com protocolo e
acompanhamento próprios. Descartamos essa ideia porque ela só muda a dor de lugar: o aluno passaria a
cuidar de dois sistemas, e a secretaria não tem motivo para adotar uma ferramenta nova.

O \textbf{Guia de Requerimentos} atua no momento anterior ao pedido. O aluno continua abrindo o
requerimento no Portal do Aluno, como hoje, mas chega lá sabendo qual categoria escolher, o que anexar,
o que escrever e quanto tempo costuma levar. Nenhum passo é acrescentado ao processo oficial.

\subsection{Funcionalidades}

\begin{enumerate}
  \item \textbf{Triagem.} O aluno descreve o que precisa com as próprias palavras (por exemplo,
        ``preciso encerrar meu estágio para começar outro''). O sistema compara a descrição com
        palavras-chave e regras cadastradas e indica a categoria exata do portal. Se não existir
        categoria adequada, ele diz isso e aponta o canal certo (Ouvidoria, coordenação, CRA).
  \item \textbf{Checklist do pedido.} Para a categoria escolhida, mostra os documentos que precisam ir
        anexados, os setores por onde o pedido vai passar e os efeitos que o aluno precisa saber antes,
        como recálculo de parcelas no trancamento ou análise da bolsa na transferência.
  \item \textbf{Modelo de texto.} Gera o texto para o campo ``Solicitação'' com as informações que o
        setor costuma pedir. O aluno copia e cola no portal.
  \item \textbf{Prazo de referência.} Alunos podem colar o ``Histórico de etapas'' de um protocolo já
        concluído. O sistema lê as datas, calcula quantos dias o pedido ficou em cada setor e guarda
        apenas o tipo de requerimento e os intervalos, sem nome, RA ou texto do pedido. Com isso,
        mostra a mediana e o percentil 75 de dias por tipo e por setor, e o aluno sabe se a espera dele
        está dentro do normal ou se já faz sentido procurar a Ouvidoria.
  \item \textbf{Registro de indeferimentos.} Quando um pedido volta negado por motivo de procedimento
        (``campo equivocado'', ``faltou documento''), o motivo vira uma regra nova da triagem, para que
        o próximo aluno não caia no mesmo erro.
\end{enumerate}

\subsection{O que o sistema não faz}

\begin{itemize}
  \item Não pede login nem senha do Portal do Aluno e não acessa o TOTVS em nome do aluno.
  \item Não envia requerimentos. O pedido oficial continua sendo feito no portal.
  \item Não depende de integração com a universidade, já que não há API pública disponível.
\end{itemize}

Essa limitação é intencional. Um sistema que pedisse a senha do portal teria problemas de segurança e
dificilmente seria aceito pela instituição.

\subsection{Arquitetura e tecnologias}

A aplicação foi feita em Java 21 com Spring Boot 4.1, seguindo a stack definida para a disciplina:
\begin{itemize}
  \item Spring Web (MVC) e Thymeleaf para as telas, com Pico.css servido pela própria aplicação e sem
        framework de front-end separado;
  \item Spring Data JPA com PostgreSQL hospedado no Neon, um serviço \emph{serverless}. Não há banco
        local para subir: cada integrante usa a própria \emph{branch} do Neon. O H2 fica só nos testes;
  \item Spring Security na área administrativa, com login por formulário. As telas do aluno continuam
        públicas e sem cadastro;
  \item JUnit 5 e Mockito para os testes, que rodam em cada \emph{pull request} pelo GitHub Actions.
\end{itemize}

O domínio aproveita bem os conceitos de orientação a objetos vistos no semestre:
\begin{itemize}
  \item \textbf{herança e polimorfismo} em \texttt{TipoRequerimento} e nas subclasses
        \texttt{RequerimentoAcademico}, \texttt{RequerimentoFinanceiro}, \texttt{RequerimentoEstagio} e
        \texttt{RequerimentoDocumental}, já que cada grupo tem alertas e um trecho próprio no modelo
        de texto;
  \item \textbf{Chain of Responsibility} na triagem (\texttt{EtapaTriagem}): a etapa das regras com
        categoria tenta reconhecer a descrição e, se não conseguir, passa para a das regras com canal
        alternativo e, por fim, para a mensagem padrão;
  \item \textbf{State} na leitura do histórico (\texttt{LeituraHistorico} e \texttt{EstadoPedido}),
        com os estados aberto, em análise em um setor, deferido e indeferido;
  \item \textbf{Strategy} no cálculo do prazo de referência (\texttt{CalculoPrazo}), com mediana e
        percentil 75, permitindo trocar um pelo outro sem mexer no resto;
  \item \textbf{Observer} na verificação da espera (\texttt{VerificadorEspera} e
        \texttt{ObservadorEspera}): cada alerta observa a situação do pedido e avisa o aluno, na
        própria tela, quando a espera passa do usual para aquele tipo.
\end{itemize}

\subsection{Quem mantém as informações}

O ponto fraco da proposta é a base de regras ficar desatualizada. A ideia é que ela seja mantida por
monitores ou pelo DCE, com uma tela simples de cadastro de categorias, documentos e regras de triagem.
Cada regra guarda a data da última revisão, e o aluno vê essa data junto da orientação.

\section{Validação}

A validação será feita com alunos da própria UCSal, comparando a abertura de requerimentos com e sem o
Guia. Não haverá envio de pedidos reais: os participantes vão simular a escolha da categoria e o
preenchimento, parando antes do botão ``Solicitar''.

\subsection{Procedimento}

\begin{enumerate}
  \item Recrutar entre 15 e 20 alunos de cursos e semestres diferentes, incluindo pelo menos alguns
        calouros e bolsistas.
  \item Aplicar um questionário curto antes do teste: quantos requerimentos já abriram, se algum foi
        indeferido por erro de categoria ou documento, quanto tempo esperaram.
  \item Dar três situações reais para cada participante, por exemplo: encerrar um estágio para começar
        outro, trancar uma disciplina por conflito com o trabalho e pedir declaração de matrícula para
        o estágio.
  \item Metade dos participantes resolve as situações só com o portal; a outra metade usa o Guia antes.
        Depois os grupos trocam de situação.
  \item Ao final, aplicar um questionário de percepção com escala de 1 a 5 e uma pergunta aberta.
\end{enumerate}

\subsection{Métricas}

\begin{table}[H]
\centering
\small
\begin{tabularx}{\textwidth}{@{}lXl@{}}
\toprule
Métrica & Como medir & Meta \\
\midrule
Acerto de categoria & categoria escolhida igual à correta, conferida pela equipe & acima de 85\% com o Guia \\
Documentos completos & anexos indicados batem com a lista exigida & acima de 80\% com o Guia \\
Tempo até o envio & cronometrado do início da situação até o ponto de clicar em ``Solicitar'' & sem piora em relação ao portal \\
Clareza percebida & média da pergunta ``eu sabia o que esperar depois de enviar'' (1 a 5) & 4 ou mais \\
\bottomrule
\end{tabularx}
\caption{Métricas de validação.}
\end{table}

O tempo até o envio é a métrica que mostra se a solução está curando a dor ou só transferindo. Se o
aluno acerta mais, mas leva o dobro do tempo, o Guia virou mais uma etapa e precisa ser simplificado.

\subsection{Resultados}

\begin{table}[H]
\centering
\small
\begin{tabular}{@{}lcc@{}}
\toprule
Métrica & Só portal & Com o Guia \\
\midrule
Acerto de categoria & --- & --- \\
Documentos completos & --- & --- \\
Tempo médio até o envio & --- & --- \\
Clareza percebida & --- & --- \\
\bottomrule
\end{tabular}
\caption{Resultados (a preencher após os testes).}
\end{table}

Além dos números, vamos levar a proposta à secretaria ou ao Registro Acadêmico para saber quais
categorias mais recebem pedidos errados. Essa conversa também serve para conferir se as regras
cadastradas estão corretas antes da demonstração.

\section*{Referências}
\raggedright

\begin{itemize}[label={},leftmargin=0pt]
  \item RECLAME AQUI. \emph{Ucsal}. Disponível em:
        \url{https://www.reclameaqui.com.br/empresa/ucsal/}. Acesso em: 21 out. 2026.
  \item RECLAME AQUI. \emph{Sem acesso ao portal do Aluno --- Ucsal}. Disponível em:
        \url{https://www.reclameaqui.com.br/ucsal/sem-acesso-ao-portal-do-aluno_a8oxIWDkHcwGYj8w/}.
        Acesso em: 21 out. 2026.
  \item RECLAME AQUI. \emph{Portal do aluno e AVA não funcionam, falta de suporte --- Ucsal}.
        Disponível em:
        \url{https://www.reclameaqui.com.br/ucsal/portal-do-aluno-e-ava-nao-funcionam-falta-de-suporte_Ilj3F59BOAl4g1Ji/}.
        Acesso em: 21 out. 2026.
  \item RECLAME AQUI. \emph{Aproveitamento de disciplinas --- Ucsal}. Disponível em:
        \url{https://www.reclameaqui.com.br/ucsal/aproveitamento-de-disciplinas_hodiNg7XaFLHC6uc/}.
        Acesso em: 21 out. 2026.
  \item RECLAME AQUI. \emph{Universidade de verdade, pra quem? --- Ucsal}. Disponível em:
        \url{https://www.reclameaqui.com.br/ucsal/universidade-de-verdade-pra-quem_16238551/}.
        Acesso em: 21 out. 2026.
  \item UNIVERSIDADE CATÓLICA DO SALVADOR. \emph{Portal do Aluno (TOTVS Educacional)}. Disponível em:
        \url{https://novoportal.ucsal.br/}. Acesso em: 21 out. 2026.
\end{itemize}

\end{document}
